% !PRINTED 382
\weberpage{382}

%% [running head, not body text -- commented out 2026-08-07] Page furniture (running
%% head, folio, signature) is excluded from this edition by standing policy.
%% \noindent 382 \hfill \textsc{Elfter Abschnitt.} \hfill \S.\,123.

Die reducirten Zahlen, deren Bilder an der Begrenzung des Grundfeldes
liegen, sind paarweise äquivalent, nämlich

\noindent a)\quad $\omega = -\tfrac{1}{2} + i\eta$ \quad und \quad $\omega_1 = \omega + 1 = \tfrac{1}{2} + i\eta$

\noindent ($a$ und $a'$ in der Figur),

\noindent b)\quad $\omega = -\xi + i\eta$ \quad und \quad $\omega_1 = -\dfrac{1}{\omega} = \xi + i\eta$

\noindent wenn $\xi^2+\eta^2 = 1$ ist ($b$ und $b'$ in der Figur). Es gilt aber ferner der
Satz:

\noindent 2.\quad \gesperrt{Von den Fällen a), b) abgesehen, sind keine zwei reducirte Zahlen äquivalent.}

Nehmen wir nämlich zwei nicht identische reducirte Zahlen an,

\weq{11}{\omega = \xi + \eta i, \quad \omega_1 = \xi_1 + \eta_1 i,}

und setzen voraus, was die Allgemeinheit nicht beeinträchtigt,

\weq{12}{\eta_1 \geqq \eta,}

so folgt aus der Formel (4):

\weq{13}{\eta_1 = \dfrac{\eta}{(\gamma\xi+\delta)^2 + \gamma^2\eta^2},}

also wegen (12):

\weq{14}{(\gamma\xi+\delta)^2 + \gamma^2\eta^2 \leqq 1,}

also auch, nach (3):

\weq{15}{1 \geqq \gamma^2\eta^2 \geqq \tfrac{3}{4}\gamma^2;}

also sind zwei Möglichkeiten:
\[
\alpha)\ \gamma = 0, \qquad \beta)\ \gamma = \pm 1.
\]

Im Falle $\alpha$) ist $\alpha\delta = 1$, also $\alpha = \delta = \pm 1$, aus (12) und (13)
ergiebt sich $\eta_1 = \eta$, ferner aus (4):
\[
\xi_1 = \xi \pm \beta.
\]
Dies ist aber nur dann mit den Bedingungen (2) verträglich, wenn
entweder $\xi_1 = \xi$, $\beta = 0$ und also $\omega_1$ mit $\omega$ identisch ist, oder wenn
$\beta = \pm 1$ und $\xi_1 = -\xi = \pm \tfrac{1}{2}$ ist, also in dem Falle a).

Im Falle $\beta$) ist nach (14) $(\xi \pm \delta)^2 \leqq 1$, und also \gesperrt{entweder}:
$\delta = 0$, $\beta\gamma = -1$ und $\xi^2+\eta^2 = 1$, $\eta_1 = \eta$, ferner nach (4)
$\xi_1 = \pm\alpha - \xi$, also, da $\xi_1 = \xi$ ausgeschlossen ist, $\xi_1 = -\xi$,
$\alpha = 0$, das ist der Ausnahmefall b),

\gesperrt{oder}: $\delta = \pm 1$, also $(\xi\pm 1)^2 + \eta^2 \leqq 1$, und folglich nach (3)
$(\xi\pm 1)^2 \leqq \tfrac{1}{4}$. Weil aber $(\xi\pm 1)$ nach (2) dem absoluten Werth nach
mindestens gleich $\tfrac{1}{2}$ sein muss, so folgt hieraus $\xi = \mp\tfrac{1}{2}$,
$\eta^2 = \tfrac{3}{4}$; mithin nach (13) $\eta_1^2 = \eta^2 = \tfrac{3}{4}$, und nach
%CONT>>
